The idea of self-organized criticality (SOC) was introduced to explain why scale-free fluctuations and burst-like dynamics appear in many driven systems without fine tuning of parameters \cite{bak1987self}. A standard example is the Bak-Tang-Wiesenfeld (BTW) sandpile, where slow driving and local relaxation lead to avalanches with a broad range of sizes \cite{Dhar1990}. Because the rules are simple but the collective behavior is nontrivial, the sandpile has become a common test case for methods that measure scaling laws and critical exponents.

Many exponent values and scaling relations for sandpile models are known from the literature, for example the study of dynamical and spatial exponents by Christensen, Fogedby, and Jensen \cite{christensen1991dynamical}. Still, in practice the numerical estimation of these exponents depends on how the system is prepared (burn-in), on finite-size effects, and on how the scaling window is chosen for the fits. In this work, we simulate several sandpile variants on finite lattices and use them to compare different driving and boundary choices and apply a clear and reproducible data analysis procedure for the exponent fits, including simple checks for reaching the stationary state.

%%%%%%%%%%%%%%%%%%%%%%%%%%%%%%%%%%%%%%%%%%%%%%%%%%%%%%%%%%%%%%%%%%%%%%%%%%%%%%%%%%%%%%%%
\subsection{Cellular Automata}
%%%%%%%%%%%%%%%%%%%%%%%%%%%%%%%%%%%%%%%%%%%%%%%%%%%%%%%%%%%%%%%%%%%%%%%%%%%%%%%%%%%%%%%%

Cellular automata (CA), first introduced by John von Neumann~\cite{vonNeumann1966SelfReproducing}, provide a discrete-time framework for modeling the evolution of spatially extended systems with local interactions. In the recent literature, e.g.\ \cite{PerezDelgadoCheung2007}, CA are often specified as a 4-tuple $(R,M,Q,\delta)$ consisting of
\begin{itemize}
    \item a lattice (cell space) $R$,
    \item a finite neighborhood $M_r \subseteq R$ for each $r\in R$,
    \item a state space $Q$,
    \item a local evolution function $\delta : Q^{\abs{M_r}} \to Q$.
\end{itemize}
A configuration at time $t$ is a map $q_t : R \to Q$ that assigns a state to each cell. The dynamics are generated by applying the same local rule at every site; the update takes the form
\begin{equation}
    q_{t+1}(r) = \delta\bigl(q_t|_{M_r}\bigr), \qquad r \in R,
\end{equation}

where $q_t|_{M_r}$ denotes the restriction of the configuration to the neighborhood of $r$.

In this work the lattice $R$ is a $d$-dimensional hypercubic grid of linear size $N$, i.e.\ $\abs{R}=N^d$, together with boundary conditions specified later. Starting from an initial configuration $q_0$, the rule $\delta$ determines the complete time evolution $\{q_t\}_{t\ge 0}$. Many physically motivated CA are additionally driven or perturbed by an external mechanism that modifies the configuration according to a separate prescription (for example, adding or subtracting values $q_t(r)\to q_t(r)\pm c$ at selected sites), after which the local update rule is applied again.

%%%%%%%%%%%%%%%%%%%%%%%%%%%%%%%%%%%%%%%%%%%%%%%%%%%%%%%%%%%%%%%%%%%%%%%%%%%%%%%%%%%%%%%%
\subsection{Self-Organized Criticality}
%%%%%%%%%%%%%%%%%%%%%%%%%%%%%%%%%%%%%%%%%%%%%%%%%%%%%%%%%%%%%%%%%%%%%%%%%%%%%%%%%%%%%%%%

Many systems in statistical physics exhibit critical behavior only when an external control parameter is tuned to a specific value. A paradigmatic example is the Ising model, which becomes critical at a particular temperature $T=T_c$ \cite{stanley1971introduction}.
At such a critical point the system becomes scale invariant meaning that correlations extend over all length scales and macroscopic observables show nontrivial scaling laws.

The system studied in this work behaves very differently. It approaches a critical state without the need to fine tune an external parameter. Instead, repeated driving and relaxation lead the system into a statistically stationary state that displays critical behavior. This stationary critical regime acts as an attractor of the dynamics and is referred to as self-organized criticality (SOC) \cite{bak1987self}.

In the SOC state, relaxation occurs through intermittent avalanches triggered by the drive. A defining feature is the absence of a characteristic avalanche scale, meaning events of many different sizes occur, from small local rearrangements up to system-spanning avalanches \cite{Dhar1990}.

%%%%%%%%%%%%%%%%%%%%%%%%%%%%%%%%%%%%%%%%%%%%%%%%%%%%%%%%%%%%%%%%%%%%%%%%%%%%%%%%%%%%%%%%
\subsection{Power Laws and Scaling} \label{sec:power-laws}
%%%%%%%%%%%%%%%%%%%%%%%%%%%%%%%%%%%%%%%%%%%%%%%%%%%%%%%%%%%%%%%%%%%%%%%%%%%%%%%%%%%%%%%%

A power law is a functional relationship in which one quantity changes proportionally to a fixed power of another. For two positive quantities $A$ and $B$, this is written as
\begin{equation}
    A(B) \propto B^{c},
\end{equation}
where $c$ is a constant exponent. Power laws are scale free. This means that rescaling the argument $B \to \lambda B$ results in a simple rescaling of the value, $A(\lambda B)=\lambda^{c}A(B)$ \cite{Newman2005}. This property makes power laws natural signatures of criticality, where no intrinsic length or time scale exists \cite{bak1987self}.

In SOC models, power laws typically arise in the distributions of avalanche observables. In a finite system, these power laws cannot extend indefinitely. The system size provides an upper cutoff \cite{christensen1991dynamical}.

Power laws are not unique to SOC and appear in many empirical contexts (for example in earthquake statistics or wealth distributions). In this work, their relevance is that they provide a quantitative way to test for scale invariance in the stationary state and to extract scaling exponents from the measured avalanche statistics.
